\section{Invariant conversion by a controlled function field}
\label{V:sec-invariant-field}

The preceding method works one cocharacter at a time and then multiplies
sections.  We now give a second nef-reduction argument which keeps
invariance throughout.  Its starting point is a sequence of invariant
twisted sections.  The base is defined by invariant functions with
pseudoeffectively controlled poles, and the final extension produces
one nonzero invariant quotient class directly.  In particular this
argument does not need the jet-generation or weight-polytope steps.

\begin{theorem}\label{V:invariant-conversion}
Let $X$ be smooth connected projective with $q(X)=0$, and let
$A=-K_X$ carry a smooth semipositive Hermitian metric.  Let an
algebraic torus $T$ act on $X$, and use the natural linearization on
$A$.  Let $G$ be a pseudoeffective $T$-linearized line bundle.
Suppose that for an unbounded set $I$ of positive integers there are
nonzero sections
\[
 \delta_j\in H^0(X,jA-G)^T\qquad(j\in I).
\]
Then $H^0(X,mA)^T\ne0$ for some integer $m>0$.
\end{theorem}

\begin{proof}
We use the maximal compact torus $T_c$ for averaging.  Average the
local weights of the given smooth metric using the natural
linearization of $A$, and call the resulting invariant smooth
semipositive metric $h_A$.  Use its pullbacks on all subsequent models.
Define $F\subset\mathbb C(X)$ to consist of the invariant rational
functions $u$ for which
\begin{equation}\label{V:controlled-poles}
 bA-D_\infty(u)\quad\hbox{is pseudoeffective for some integer }b>0,
\end{equation}
and include zero.  Here $D_\infty(u)$ denotes the effective polar
divisor on $X$.  Sum and product preserve this condition by the
usual pole inequalities, and inversion preserves it because a
rational function's zero and pole divisors are linearly equivalent.
Thus $F$ is a field.

It is relatively algebraically closed in $\mathbb C(X)$.  For if
$u^d+\sum_{i<d}c_i u^i=0$ with $c_i\in F$, then at a pole of $u$
of order $a>0$ some term satisfies
$\operatorname{ord}(D_\infty(c_i))\ge(d-i)a$.  Consequently
$D_\infty(u)\le\sum_{c_i\ne0}D_\infty(c_i)$.
The root is invariant as well: the connected torus permutes the
finite set of roots trivially.  It therefore belongs to $F$.
The field $F$ is finitely generated over $\mathbb C$.  Indeed take
a transcendence basis $x$ for $F/\mathbb C$ and extend it by $z$ to
one for $\mathbb C(X)/\mathbb C$.  The finite degree of
$\mathbb C(X)/\mathbb C(x,z)$ bounds $[F:\mathbb C(x)]$, because
$z$ is algebraically independent over $F$.  Hence the latter degree
is finite.

Choose a projective model of $F$ and apply
Lemma~\ref{V:flat-model} equivariantly.  Denote the smooth base by
$Y$, the normal equidimensional model by $V$, and a smooth equivariant
resolution by $S$, with maps
\[
 p:V\to X,\quad f:V\to Y,
 \qquad p:S\to X,\quad f:S\to Y.
\]
The base action is trivial and $f$ has connected fibers.  Put
$A_S=p^*A$ and $E=K_S-p^*K_X\ge0$.  The natural section
$s_E$ is invariant; this follows directly from pullback of top forms.

We record a domination consequence of the definition of $F$.
For every smooth birational model of $X$ mapping to a smooth
projective base $Y_1$ with $\mathbb C(Y_1)\subset F$, and every ample
Cartier divisor $H$ on $Y_1$, there is an integer $b>0$ with
\begin{equation}\label{V:field-domination}
 b p^*A-f_1^*H\quad\hbox{pseudoeffective}.
\end{equation}
Choose sections $s_0,\ldots,s_N$ generating $hH$ with $s_0\ne0$.
The ratios $s_i/s_0$ are in $F$.  On every prime of the total space,
generation ensures that one pulled-back section has zero order, so
\[
 \operatorname{div}(f_1^*s_0)
 \le p^*\sum_iD_\infty(s_i/s_0).
\]
The pole inequality is valid also on exceptional primes, since the
pole divisor of the pullback of a rational function is bounded by
the pullback of its polar divisor.  Sum the pseudoeffective bounds
\eqref{V:controlled-poles} and then add $(h-1)f_1^*H$ to obtain
\eqref{V:field-domination}.

\subsection*{An invariant relative section with vertical minima}

Choose $r<s$ in $I$.  For $j\in I$ with $j>s$, the invariant function
\[
 u_j=\frac{\delta_j^{s-r}\delta_r^{j-s}}
             {\delta_s^{j-r}}
\]
belongs to $F$.  Its poles are bounded by
$(j-r)\operatorname{div}(\delta_s)$, and
$sA-\operatorname{div}(\delta_s)\sim G$ is pseudoeffective.
At each horizontal prime of $V\to Y$, the order of $u_j$ is zero.
If $v_j$ denotes the order there of $p^*\delta_j$, then
\[
 (s-r)v_j=(j-r)v_s-(j-s)v_r.
\]
Since $v_j\ge0$ for arbitrarily large $j$, we have $v_s\ge v_r$.
Thus the invariant rational section
$t=p^*(\delta_s/\delta_r)$ of $(s-r)p^*A$ has no horizontal pole.
If $Y$ is a point, $t$ is regular by normality and descends to $X$,
proving the assertion.

Otherwise, over each prime $Q$ of $Y$, take the minimum of
$\operatorname{ord}_\Gamma(t)/\operatorname{ord}_\Gamma(f^*Q)$
over components $\Gamma$ dominating $Q$.  Let $k_0$ clear their
finitely many nonzero denominators.  Subtracting the pullback of the
resulting Cartier divisor $B$ on $Y$ gives an invariant section
\begin{equation}\label{V:relative-invariant}
 0\ne\tau\in H^0(V,a p^*A-f^*B)^T,
 \qquad a=k_0(s-r)>0.
\end{equation}
It has zero order on at least one component over every base prime.
There are no vertical primes above higher codimension on $V$, so
nonnegative orders on these primes and on horizontal primes prove
regularity.  Pull $\tau$ to $S$, and write $P_\tau$ for its effective
divisor.  Every base line bundle in this construction has trivial
linearization.

\subsection*{Rank one for the invariant direct image}

For each $l\ge1$ consider
\[
 \mathcal E_l=f_*(K_{S/Y}+(al+1)A_S).
\]
Over a dense base-change open set its invariant part has rank one.
To prove the upper bound, twist by a sufficiently ample $H$ on $Y$
so that $H-K_Y$ is ample and the twisted direct image is generated.
Averaging shows that invariant global sections generate the invariant
fiber spaces there.  The ratio of any two invariant global sections
lies in $F$: if the denominator has divisor $N$ on $S$, then
\[
 N\sim alA_S+E+f^*(H-K_Y).
\]
Equation~\eqref{V:field-domination} bounds the last summand
pseudoeffectively by a multiple of $A_S$.  Pushforward to the smooth
$X$ preserves pseudoeffectivity and kills $E$, so the polar divisor
of the ratio on $X$ satisfies \eqref{V:controlled-poles}.
Since all these ratios come from $Y$, the invariant fiber space is
at most one-dimensional.  It is at least one-dimensional, because
in local frames $\beta$ of $B$ and $\tau_Y$ of $-K_Y$ the section
\[
 v_l=(\tau f^*\beta)^l s_E f^*\tau_Y
\]
is a nonzero invariant relative canonical section with coefficient
$(al+1)A_S$.

We can now apply the invariant version of
Lemma~\ref{V:bounded-maximum}.  The action on the base and its line
$B$ is trivial, $s_E$ is the natural invariant Jacobian section,
$\tau$ is invariant, and the coefficient metric is smooth and
compact-invariant.  The rank-one assertion has been proved for
\emph{every} $l\ge1$.  The lemma therefore gives locally bounded
semipositive weights on $B$, so $B$ is nef.  Notice that this step
uses the invariant adjoint rank and the orthogonal averaging
projection; ordinary rank one has not been assumed.

\subsection*{Nef reduction and a nonzero quotient class}

Since $Y$ is dominated rationally by $X$, pullback of holomorphic
one-forms gives $q(Y)=0$.  Apply Lemma~\ref{V:nef-descent}, lift $S$
further equivariantly, and clear the common denominators.  Retaining
$p:S\to X$, $A_S=p^*A$, $E=K_S-p^*K_X$, and writing $g:S\to Z$,
we have a nef Cartier divisor $L$ and an invariant section
\[
 \tau\in H^0(S,aA_S-g^*L)^T,
 \qquad P_\tau=\operatorname{div}(\tau)\ge0.
\]
The line-bundle identification comes from a base with trivial action,
so invariance is preserved.  A zero-dimensional $Z$ finishes the
proof by descending $\tau$.  Otherwise let $b=\dim Z>0$.

Fix very ample $H$ on $Z$.  Equation~\eqref{V:field-domination},
using $\mathbb C(Z)\subset F$, gives a pseudoeffective divisor
$R_0=b_0A_S-g^*H$ for some integer $b_0>0$.
Take a $T_c$-invariant semipositive singular metric $h_R$ on it.
Such a metric is obtained by averaging its local weights relative
to an invariant smooth reference metric.  The quasi-plurisubharmonic
potentials have uniformly bounded $L^1$ norm under the compact group,
so the average is well defined, has semipositive curvature
distributionally, and gives invariant plurisubharmonic weights.
Choose a small rational $\epsilon>0$ with
$\mathcal I(h_R^\epsilon)=\mathcal O_S$.

Apply Lemma~\ref{V:sesha} with $c=2b$.  For sufficiently large $k$,
the ample rational divisor $B_k=kL+\epsilon H$ has Seshadri constant
at least $4b$ at very general points.  Choose such a point $z$ with
some $x\in g^{-1}(z)$ where $g$ is submersive and $P_\tau\cup E$
is absent.  Put $\beta=3b$.  The pole construction in the same lemma,
applied to $B_k/2$ with $c=\beta/2$, gives a smooth positive metric
of curvature $\omega_k/2$ and a function $\psi$ with neat analytic
singularities, smooth away from $z$, such that
\[
 i\partial\bar\partial\psi\ge-\tfrac12\omega_k,
 \qquad \psi\le\tfrac{\beta}{2}\log|w|^2+O(1).
\]
Give $B_k$ the induced smooth metric of curvature $\omega_k$, and
put $\Psi=g^*\psi$.  As in the first-jet argument, surjectivity of
$g$ preserves the local analytic-singularity form and prevents the
pullback from being identically $-\infty$.  Thus $\Psi$ has neat
analytic singularities, is smooth away from $g^{-1}(z)$, and satisfies
\[
 i\partial\bar\partial\Psi\ge-\tfrac12 g^*\omega_k,
 \qquad
 \Psi\le\tfrac{\beta}{2}\log|w|^2+O(1)
\]
in transverse base coordinates at the chosen submersion point.

Choose an integer $p_0\ge k$ sufficiently large, set $m=ap_0$, and
use the decomposition
\begin{equation}\label{V:invariant-metric-decomposition}
 (m+1)A_S
 =k(aA_S-g^*L)+g^*B_k+\epsilon R_0
        +(m+1-ka-\epsilon b_0)A_S.
\end{equation}
The last coefficient is nonnegative.  Give the four terms respectively
the divisor metric of $\tau^k$, the indicated smooth base metric,
$h_R^\epsilon$, and the remaining smooth semipositive metric.
They define a $T_c$-invariant metric $h_J$ on $J=(m+1)A_S$ with
\[
 i\Theta_{h_J}+(1+s)i\partial\bar\partial\Psi\ge0
 \qquad(0\le s\le1).
\]
Its multiplier ideal is exactly
\begin{equation}\label{V:fixed-divisor-ideal}
 \mathcal I(h_J)=\mathcal O_S(-kP_\tau).
\end{equation}
Divisibility by the equation of $kP_\tau$ is sufficient because the
remaining singular factor $h_R^\epsilon$ is integrable.  It is
necessary at the generic smooth point of each prime divisor: the
remaining local weights are bounded above, and an uncancelled
integral pole already prevents square integrability there.  The
resulting codimension-one divisibility implies divisibility as a
holomorphic germ, since the local rings of the smooth variety are
factorial.  This proves \eqref{V:fixed-divisor-ideal}.

Apply \cite[Theorem~1.1]{CDM2017}.  It gives
\[
 H^0\!\left(S,\mathcal O_S(K_S+J)\otimes\mathcal I(h_J)\right)
 \longrightarrow
 H^0\left(S,\mathcal O_S(K_S+J)\otimes
       \frac{\mathcal I(h_J)}{\mathcal I(h_Je^{-\Psi})}\right)
 \quad\hbox{surjectively}.
\]
The quotient is supported over $z$.  In a local frame $e_L$ of $L$
near $z$, the section
\[
 s_E(\tau g^*e_L)^{p_0}
\]
is a section of the larger ideal twist near the whole fiber, because
$p_0\ge k$.  Its quotient class glues with zero away from that fiber
and is invariant, as the base action and its frame action are trivial.
The class is nonzero.  At $x$ its coefficient is a unit; the local
weight of $h_J$ is bounded above, while the transverse pole has
exponent $\beta/2=3b/2\ge b$.  Thus its squared norm for
$h_Je^{-\Psi}$ is not locally integrable in the $b$ transverse
variables.

Lift this nonzero class and average the lift over $T_c$.
The ideals are invariant, the map is equivariant, and averaging
preserves the chosen invariant quotient class.  We obtain a nonzero
invariant section of
\[
 K_S+J=mA_S+E.
\]
Remove the canonical exceptional factor and push to $X$.
It is a section of $mA$ off a codimension-two subset, hence extends
holomorphically over the smooth $X$.  It remains nonzero and
$T_c$-invariant, therefore $T$-invariant.  This proves the theorem.
\end{proof}

\begin{corollary}\label{V:invariant-differentials}
Let $X,T,A$ satisfy the geometric hypotheses of
Theorem~\ref{V:invariant-conversion}.  If for a fixed
$0\le p\le\dim X$ there are nonzero invariant sections of
$\Omega_X^p\otimes jA$ for an unbounded set of positive integers
$j$, then $H^0(X,mA)^T\ne0$ for some $m>0$.
\end{corollary}

\begin{proof}
For $p=0$ the conclusion is part of the assumption.  Otherwise use
the determinant construction of Section~\ref{V:sec-extremal}, now
with invariant forms throughout.  Form
the saturated invariant subsheaf $\mathcal S\subset\Omega_X^p$
spanned generically by the chosen forms, and put
$G=-(\bigwedge^{\operatorname{rank}\mathcal S}\mathcal S)^{**}$.
This is a line bundle on the smooth $X$, with its induced
linearization.  The inclusion into the tensor power of $\Omega_X^1$
and \cite[Theorem~4.1]{LMPTX2023} make $G$ pseudoeffective.
A form with arbitrarily large twist can be completed to a generic
basis using a fixed finite collection of forms.  Their wedges,
followed by the reflexive determinant map, give nonzero invariant
sections of $j'A-G$, with $j'$ unbounded.  Theorem~\ref{V:invariant-conversion}
applies.
\end{proof}

\begin{remark}\label{V:invariant-index-application}
Under the geometric hypotheses of Theorem~\ref{V:invariant-conversion},
in particular $q(X)=0$, this conversion also supplies the nef-reduction
route from invariant cohomology.  The invariant-index theorem
\cite[Theorem \textup{(Invariant anticanonical index)}]{CompanionI} applies to the compact
complex manifold $X$ and its maximal compact torus, with the natural
linearization of $A=-K_X$.  It says that the invariant index agrees
with a rational polynomial on every positive multiple of some
integer, and that this polynomial's value at zero is the actual
invariant index at zero.  The connected torus acts trivially on
ordinary cohomology by homotopy and therefore on its Hodge summands.
Consequently this zero value is $\chi(X,\mathcal O_X)$.  If it is nonzero,
some fixed $q$ has nonzero invariant $H^q(X,mA)$ for unbounded $m$.
The compact-averaged hard Lefschetz map with coefficient $(m+1)A$
produces the invariant twisted forms needed by
Corollary~\ref{V:invariant-differentials}.  Thus
$H^0(X,mA)^T\ne0$ for some $m>0$ whenever
$\chi(X,\mathcal O_X)\ne0$.  In particular the hypothesis $H^q(X,\mathcal O_X)=0$ for every $q>0$ gives the zero
index value $1$ and $q(X)=0$.  The index theorem and the conversion
play separate roles: the former retains the actual constant term,
whereas the latter produces an invariant section without presuming
an equivariant version of an ordinary nonvanishing criterion.
\end{remark}

The compact-factor application of these invariant conclusions is
given in the descent section.  Only the geometric structure and
finite-cover norm are needed there; neither of the proofs above
uses the headline nonvanishing theorem.
